\section{Search Strategy}
\label{sec:search}

The model adopts a constructive search strategy that exploits a structural
property of the problem: the pallet can be filled row by row, top to
bottom, and within each row left to right, taking the longest pick the
gripper allows at each step. Under this plan, the row above each new chunk
is already full when it is placed, so the collision rule naturally forces
the plate to open downward into the still-empty rows below, where no
collision is possible. The plan is feasible by construction, and a search
that follows this order finds a solution with essentially no backtracking.

For each chunk the search must fix four values together: the pallet row
\texttt{y\_row[i]}, the length \texttt{chunk\_len[i]}, the starting column
\texttt{x\_start[i]}, and the plate direction \texttt{plate[i]}. The
desired values are, respectively, the smallest available row, the longest
pick possible, the leftmost column, and the downward plate.

This intent cannot be expressed directly. The \texttt{int\_search}
annotation visits a single array in index order and applies a single
value-selection rule, whereas the strategy above needs four arrays at once
and mixed extrema: \texttt{indomain\_min} on row and column, but
\texttt{indomain\_max} on length and plate. The model works around both
limitations with a single construction: it channels the four decisions of
each chunk into one \texttt{branch} array, laid out in blocks of four, and
applies a uniform \texttt{indomain\_min} to the whole array. For the
variables on which a maximum is wanted, the entry stored in
\texttt{branch} is the \emph{complement}, so that minimizing the entry
yields the maximum of the original variable.

\vspace{0.5cm}
\begin{lstlisting}
par int: branch_ub = max([boxes_along_y, boxes_along_x, max_pkble_boxes]);

array[1..4*n_chunks] of var 0..branch_ub: branch;

constraint forall(i in 1..n_chunks)(
  branch[4*i-3] = y_row[i]                       /\
  branch[4*i-2] = max_pkble_boxes - chunk_len[i] /\
  branch[4*i-1] = x_start[i]                     /\
  branch[4*i]   = 1 - plate[i]
);

solve :: int_search(branch, input_order, indomain_min)
      satisfy;
\end{lstlisting}
\vspace{0.5cm}

The branch array is laid out so that entries \(4i-3\), \(4i-2\),
\(4i-1\), \(4i\) correspond to the row, length, column, and plate of
chunk \(i\). Visiting the array in \texttt{input\_order} therefore
processes all four decisions of chunk \(1\) before moving to chunk
\(2\), and so on.

For row and column, the entries store \texttt{y\_row[i]} and
\texttt{x\_start[i]} directly, so \texttt{indomain\_min} selects the
smallest available values: the topmost row and the leftmost column. For
length and plate the desired behavior is the opposite: the search should
prefer the longest pick and the downward plate. Since \texttt{indomain\_min}
always minimizes, these two entries store the \emph{complement} of the
original variable. The entry \texttt{max\_pkble\_boxes - chunk\_len[i]}
reaches its minimum of \(0\) when \texttt{chunk\_len[i] = max\_pkble\_boxes};
the entry \texttt{1 - plate[i]} reaches its minimum of \(0\) when
\texttt{plate[i] = 1}.

Figure~\ref{fig:sample-schedule} shows an example output of the proposed model.

\begin{figure}[h]
  \centering

  \begin{subfigure}[t]{0.33\linewidth}
    \centering
    \includegraphics[width=\linewidth]{figures/schedule/schedule_step_01.pdf}
    \caption{Step 1}
  \end{subfigure}\hfill
  \begin{subfigure}[t]{0.33\linewidth}
    \centering
    \includegraphics[width=\linewidth]{figures/schedule/schedule_step_03.pdf}
    \caption{Step 3}
  \end{subfigure}\hfill
  \begin{subfigure}[t]{0.33\linewidth}
    \centering
    \includegraphics[width=\linewidth]{figures/schedule/schedule_step_05.pdf}
    \caption{Step 5}
  \end{subfigure}

  \vspace{0.3cm}

  \begin{subfigure}[t]{0.33\linewidth}
    \centering
    \includegraphics[width=\linewidth]{figures/schedule/schedule_step_07.pdf}
    \caption{Step 7}
  \end{subfigure}\hfill
  \begin{subfigure}[t]{0.33\linewidth}
    \centering
    \includegraphics[width=\linewidth]{figures/schedule/schedule_step_09.pdf}
    \caption{Step 9}
  \end{subfigure}\hfill
  \begin{subfigure}[t]{0.33\linewidth}
    \centering
    \includegraphics[width=\linewidth]{figures/schedule/schedule_step_11.pdf}
    \caption{Step 11}
  \end{subfigure}

  \vspace{0.3cm}

  \begin{subfigure}[t]{0.33\linewidth}
    \centering
    \includegraphics[width=\linewidth]{figures/schedule/schedule_step_13.pdf}
    \caption{Step 13}
  \end{subfigure}\hfill
  \begin{subfigure}[t]{0.33\linewidth}
    \centering
    \includegraphics[width=\linewidth]{figures/schedule/schedule_step_14.pdf}
    \caption{Step 14}
  \end{subfigure}\hfill
  \begin{subfigure}[t]{0.33\linewidth}
    \centering
    \includegraphics[width=\linewidth]{figures/schedule/schedule_step_15.pdf}
    \caption{Step 15}
  \end{subfigure}

  \caption{Example of a feasible sequence. The red border marks the
  chunk being placed at the selected step, while the blue arrow indicates the
  opening side of the plate.}
  \label{fig:sample-schedule}
\end{figure}